\documentclass[11pt,a4paper]{article}
\usepackage[T1]{fontenc}
\usepackage[utf8]{inputenc}
\usepackage{lmodern}
\usepackage[margin=25mm,headheight=15pt]{geometry}
\usepackage{amsmath,amssymb,amsthm,mathtools,mathrsfs}
\usepackage{microtype,booktabs,array}
\usepackage{xcolor,fancyhdr}
\usepackage{hyperref}
\usepackage{bookmark}
\definecolor{Ink}{HTML}{253640}
\definecolor{Accent}{HTML}{536044}
\hypersetup{colorlinks=true,linkcolor=Ink,citecolor=Accent,urlcolor=Ink,
 pdftitle={Infinitesimal Resonance and Hidden Monodromy in Surcomplex Dynamics},
 pdfauthor={Research article prepared for Vladimir Reshetnikov},
 pdfsubject={Hahn-algebraicity, common-domain linearization, and sharp infinitesimal scales}}
\pagestyle{fancy}
\fancyhf{}
\fancyhead[L]{\small Infinitesimal resonance and hidden monodromy}
\fancyhead[R]{\small Surcomplex dynamics}
\fancyfoot[C]{\thepage}
\setlength{\emergencystretch}{2em}
\numberwithin{equation}{section}
\newtheorem{theorem}{Theorem}[section]
\newtheorem{proposition}[theorem]{Proposition}
\newtheorem{lemma}[theorem]{Lemma}
\newtheorem{corollary}[theorem]{Corollary}
\theoremstyle{definition}
\newtheorem{definition}[theorem]{Definition}
\newtheorem{example}[theorem]{Example}
\theoremstyle{remark}
\newtheorem{remark}[theorem]{Remark}
\newcommand{\C}{\mathbb C}
\newcommand{\R}{\mathbb R}
\newcommand{\Q}{\mathbb Q}
\newcommand{\N}{\mathbb N}
\newcommand{\Z}{\mathbb Z}
\newcommand{\No}{\mathbf{No}}
\newcommand{\K}{K_\Gamma}
\newcommand{\A}{\mathscr A_\Gamma}
\newcommand{\B}{\mathscr B_\Gamma}
\newcommand{\m}{\mathfrak m_\Gamma}
\newcommand{\eps}{\varepsilon}
\newcommand{\dd}{\mathrm d}
\newcommand{\id}{\operatorname{id}}
\newcommand{\supp}{\operatorname{supp}}
\newcommand{\Hol}{\mathcal O}
\newcommand{\Mer}{\mathcal M}
\newcommand{\Res}{\operatorname{Res}}
\newcommand{\Log}{\operatorname{Log}}
\newcommand{\ExpH}{\operatorname{Exp}_{\!H}}
\newcommand{\Neu}[1]{\langle #1\rangle}
\newcommand{\coef}[2]{[#1]#2}
\newcommand{\val}{v}
\newcommand{\ord}{\operatorname{ord}}
\newcommand{\Ucov}{\widetilde X}
\newcommand{\commit}{\texttt{39f2be6667ade51bca2b45daa47e289d69c09764}}
\newcommand{\smallO}[1]{O(#1)}
\DeclareMathOperator{\lc}{lc}

\title{\textbf{Infinitesimal Resonance and Hidden Monodromy\\
in Surcomplex Dynamics}\\[7pt]
\large A Hahn-algebraicity dichotomy, common-domain linearization,\\
and sharp infinitesimal scales}
\author{Research article prepared for Vladimir Reshetnikov}
\date{21 September 2026}

\begin{document}
\maketitle
\begin{abstract}
Let $\eps=t^\delta$ be a positive-valuation Hahn monomial in a set-sized
surcomplex workspace, and let $V\in\C[w]$ satisfy $V(0)=0$ and
$V'(0)=\mu\ne0$.  We construct the normalized Koenigs coordinate of
$G_\eps(w)=w+\eps V(w)$ as a Hahn series whose coefficients are holomorphic
on one common ordinary domain.  We prove an algebraicity dichotomy:
this coordinate is algebraic over the Hahn field
$\C(w)((t^\Gamma))$ if and only if $V$ is linear.  For a simple zero $a$ of
$V$, its exact coefficientwise monodromy is
\[
 \ExpH\!\left(2\pi i\,
 \frac{\Log(1+\mu\eps)}{\Log(1+V'(a)\eps)}\right).
\]
An infinitesimal part in this exponent cannot be removed by any finite
branched cover.  This explains how an algebraic leading coordinate can
acquire infinite-order monodromy at its first infinitesimal correction.
For $f(z)=(1+\eps)z+z^{q+1}$ we also determine the exact strong Taylor
summability domain, $v(z)>\delta/q$, prove that the conjugacy there is a
valuative isometry, and construct a coefficientwise continuation at the
boundary scale.  Full proofs separate ordinary analytic continuation from
Hahn summability and work in arbitrary-rank value groups.  A standard-library
Python program supplies 2,764 exact finite consistency checks.  The
algebraicity and monodromy statements are proposed research contributions;
no priority certification or resolution of a named published conjecture is
claimed.
\end{abstract}

\medskip
\noindent\textbf{Research status.}
This is an unrefereed, AI-generated research draft.  The arguments below are
mathematical proofs offered for scrutiny, not proofs checked by a proof
assistant.  The computational appendix tests finite identities and does not
replace the proofs.

\clearpage
\setcounter{tocdepth}{2}
{\small\tableofcontents}
\clearpage

\section{The result and why it is not ordinary complex linearization}
\label{sec:intro}

A small change in a multiplier can be invisible in the leading complex
geometry and decisive in the full surcomplex geometry.  Our basic example is
\begin{equation}\label{eq:basicexample}
 G_\eps(w)=w+\eps(w+w^2).
\end{equation}
Its leading normalized linearizing coordinate is the rational function
$w/(1+w)$.  The next coefficient is
\[
 \frac{w}{1+w}\log(1+w).
\]
Thus a rational leading coordinate does not even predict the branching type
of the first correction.  The full answer has infinite-order monodromy
around $w=-1$, although the leading answer has none.  This is not a numerical
small-divisor estimate: the monodromy has a genuinely infinitesimal component.

Throughout, $\Gamma$ is a divisible, set-sized ordered abelian subgroup of
$\No$, $t^\gamma=\omega^{-\gamma}$ denotes the Conway monomial, and
\[
 \K=\C((t^\Gamma))\subset\No[i].
\]
All supports are well-ordered sets, and every sum is subject to the precise
strong summability rule in Section~\ref{sec:hahn}.  We fix $\delta>0$ and
write $\eps=t^\delta$.  Ordinary powers of $\eps$ therefore form a particularly
simple Hahn subring $\C[[\eps]]$.  In particular, $\eps$ is not a small
nonzero complex number.

\subsection{The principal theorem}

For a polynomial $V$ with $V(0)=0$ and $V'(0)=\mu\ne0$, a normalized
Koenigs coordinate means a coefficientwise holomorphic solution of
\begin{equation}\label{eq:schroederintro}
 H_\eps\bigl(w+\eps V(w)\bigr)
   =(1+\mu\eps)H_\eps(w),\qquad
 H_\eps(0)=0,\quad H_\eps'(0)=1.
\end{equation}
Composition here is Taylor substitution in the positive Hahn displacement
$\eps V(w)$.  It is defined without assuming that the ordinary domain is
mapped into itself by a numerically specialized polynomial.

\begin{theorem}[Polynomial Euler maps: algebraicity and monodromy]
\label{thm:main}
Let $V\in\C[w]$ satisfy $V(0)=0$ and $V'(0)=\mu\ne0$.
On every simply connected ordinary domain $U$ containing $0$ and no other
zero of $V$, equation~\eqref{eq:schroederintro} has a unique solution
\[
 H_\eps\in\Hol(U)[[\eps]],
 \qquad H_\eps(w)=\sum_{n\ge0}\eps^n h_n(w).
\]
Its coefficients admit analytic continuation to the universal cover of
$X=\C\setminus Z(V)$.  In the corresponding field of meromorphic-coefficient
Hahn series:
\begin{equation}\label{eq:dichotomyintro}
 H_\eps\text{ is algebraic over }\B:=\C(w)((t^\Gamma))
 \quad\Longleftrightarrow\quad \deg V=1.
\end{equation}
At a simple zero $a$ of $V$, put $\nu_a=V'(a)$.  Positive continuation once
around $a$ multiplies $H_\eps$ by
\begin{equation}\label{eq:monodromyintro}
 M_a=\ExpH(2\pi i\rho_a),\qquad
 \rho_a=\frac{\Log(1+\mu\eps)}{\Log(1+\nu_a\eps)}.
\end{equation}
This multiplier has finite order precisely when $\nu_a=\mu$; in that case it
is $1$.  If $V$ is nonlinear and all its zeros are simple, at least one zero
has infinite-order monodromy, which no finite branched cover removes.
\end{theorem}

The field in~\eqref{eq:dichotomyintro} is deliberately large.  It consists of
Hahn series with rational-function coefficients; it is \emph{not} merely
the rational-function field $\K(w)$.  In particular, proving transcendence
over $\B$ is stronger than proving transcendence over $\K(w)$.  The word
``transcendence'' concerns a \emph{function} in an ambient meromorphic Hahn
field.  It says nothing contradictory about the algebraic closedness of
$\No[i]$, and it is not a claim about the value $H_\eps(w_0)$ at an individual
surcomplex point.

The proof has three distinct parts.  A strong logarithm of substitution
produces a holomorphic vector field, whose ordinary coefficientwise
integration gives the coordinate.  The derivative of that vector field at a
fixed point is the Hahn logarithm of the corresponding multiplier.  Finally,
monodromy gives an infinite orbit under automorphisms fixing $\B$, whereas an
algebraic element can have only a finite orbit.  Multiple zeros of $V$ are
handled separately by the already transcendental leading coefficient.

\subsection{A sharp resonant model}

The rescaling $z=s w$, with $s=t^{\delta/q}$ and $s^q=\eps$, changes
\begin{equation}\label{eq:originalmodelintro}
 f_{q,\eps}(z)=(1+\eps)z+z^{q+1}
\end{equation}
into $w+\eps(w+w^{q+1})$.  Let $h$ be the normalized formal linearizer of
$f_{q,\eps}$.  We will prove that its Taylor family is strongly summable
\emph{exactly} at
\begin{equation}\label{eq:sharpdomainintro}
 D_{q,\delta}=\{z\in\K:v(z)>\delta/q\}.
\end{equation}
On this disk it is an invertible conjugacy and satisfies
$v(h(x)-h(y))=v(x-y)$.  Nevertheless, after rescaling, its coefficients are
holomorphic on $|w|<1$ and can be continued around the nonzero fixed points
$w^q=-1$.  Their monodromy exponent is
\[
 \rho_q(\eps)=\frac{\Log(1+\eps)}{\Log(1-q\eps)}.
\]
The leading coordinate $w(1+w^q)^{-1/q}$ becomes single-valued after finite
ramification.  The full coordinate never does.  This is the interaction
between a sharp infinitesimal scale and an obstruction invisible in leading
complex geometry.

\subsection{Relation to the repository and the literature}
\label{sec:related}

The repository was inspected at commit \commit\,\cite{Repo}.  Its documentation
map organizes surcomplex analysis into foundational analysis, analytic
geometry, finite deformations, contours, global divisors, polynomial algebra,
and trigonometry.  The present article adds local dynamical conjugacy and
coefficientwise monodromy to that documented program; it does not re-prove
its divisor or residue-duality results.  In particular, we adopt its crucial
distinction between a common ordinary domain for all coefficient functions
and a family of germs with no common radius.  We do not rely on the
unrefereed advanced results of those reports: the arguments needed here are
given below.

The underlying mechanisms have substantial precedents.  Neumann's support
lemma is the classical algebraic tool behind generalized series
\cite{Neumann}.  Formal embedding of near-identity maps in autonomous flows
is well established; Gelfreich and Vieiro discuss that formal construction
and the difference between formal interpolation and analytic embedding at a
numerical parameter\,\cite{GV}.  Strong derivations and strong automorphisms
also occur in current work on surreal Hahn fields\,\cite{KKS}.
Our logarithm-of-substitution lemma is an explicit common-domain version of
that machinery, not a claim to have invented formal logarithms.

Likewise, equal-characteristic-zero ultrametric linearization, including
sharp estimates and examples where nearby periodic points bound a
linearization disk, is treated by Lindahl\,\cite{Lindahl}.  The quadratic
scale in our model belongs to that established landscape.  What is pursued
here is a different combination: arbitrary-rank strong support control,
continuation of ordinary coefficient functions at the critical scale, an
exact infinitesimal monodromy formula, and transcendence even over the
rational-coefficient Hahn field.  We do not present a disk estimate alone as
a new theorem of classical dynamics.

\paragraph{Novelty boundary.}
A targeted literature check located these adjacent theories but no matching
statement of Theorem~\ref{thm:main} over $\C(w)((t^\Gamma))$, or of the explicit
finite-cover obstruction developed here.  These are therefore proposed
contributions, to the author's present knowledge, not certified priority
claims.  This article takes the ``prove new theorems'' alternative in the
request; it does not advertise the solution of a named published open
problem.  The limitations of the construction are collected in
Section~\ref{sec:limits}.

\section{Hahn analysis on one common ordinary domain}
\label{sec:hahn}

\subsection{Fields, supports, and the relevant notion of sum}

A Hahn series over an ordinary coefficient field $k$ is an expression
\[
 x=\sum_{\gamma\in S}x_\gamma t^\gamma,
 \qquad x_\gamma\in k,
\]
where $S\subset\Gamma$ is well ordered.  Its valuation is the least exponent
with nonzero coefficient, and $v(0)=+\infty$.  Multiplication is convolution.
The valuation is trivial on nonzero constants in $\C$.  It consequently
ignores the ordinary numerical size of a complex coefficient: a very large
complex derivative and a very small nonzero complex derivative both have
Hahn valuation zero.

\begin{definition}[Strong summability]
A set-indexed family $(x_j)_{j\in J}$ is \emph{strongly summable} when
$\bigcup_j\supp(x_j)$ is well ordered and each exponent occurs in the support
of only finitely many $x_j$.  Its sum is defined coefficientwise.  We use the
same definition with coefficients in an ordinary commutative ring.
\end{definition}

This definition does not permit cancellation to rescue an inadmissible
family.  Infinitely many nonzero terms at a single exponent already violate
summability, even when an unrelated summation method might assign a value to
their complex coefficients.

\begin{lemma}[Neumann support lemma]\label{lem:neumann}
Let $S\subset\Gamma_{>0}$ be well ordered.  The additive monoid
$\Neu S$ generated by $S$, including $0$, is well ordered.  Every exponent
has only finitely many representations as a finite word with letters in
$S$.  If $T\subset\Gamma$ is well ordered, then $T+\Neu S$ is well ordered,
and a fixed exponent has only finitely many representations as an element
of $T$ plus such a word.
\end{lemma}

This is the standard generalized-series lemma\,\cite{Neumann}; the
finite-word version is equivalent to the usual finite-factorization version
because each fixed finite word has finitely many permutations.  The last
assertion combines that lemma with the usual finite-decomposition property
of two well-ordered subsets of an ordered abelian group.  Indeed, infinitely
many distinct decompositions of one exponent would force one of the two
components into an infinite strictly descending sequence.

All infinite operator manipulations below are certified by this lemma.
For example, if $v(u)>0$, then
\begin{equation}\label{eq:logexp}
 \Log(1+u)=\sum_{n\ge1}\frac{(-1)^{n+1}}n u^n,
 \qquad
 \ExpH(u)=\sum_{n\ge0}\frac{u^n}{n!}
\end{equation}
are strong sums.  The proof is not the assertion that $n v(u)$ eventually
exceeds every element of $\Gamma$.  That assertion is false in higher-rank
ordered groups, and is nowhere required.

For a finite Hahn constant $c=c_0+c_+$ with $c_0\in\C$ and $v(c_+)>0$, set
\begin{equation}\label{eq:finiteexp}
 \ExpH(c)=e^{c_0}\sum_{n\ge0}\frac{c_+^n}{n!}.
\end{equation}
Finite-word regrouping proves the usual exponential product law.  In
particular $\ExpH(\Log(1+u))=1+u$ for positive $u$.

\subsection{Coherent coefficient rings}

Fix an ordinary connected open set $U\subset\C$.  We work in
\[
 \A(U)=\Hol(U)((t^\Gamma)).
\]
An element has a single well-ordered exponent support, and every coefficient
is holomorphic on this same $U$.  Its derivative is taken coefficientwise.
A primitive, when it exists, is also taken coefficientwise; no estimate
uniform in the exponent is part of the definition.

If a series has a nowhere-zero leading coefficient on $U$, it is invertible
in $\A(U)$.  To see this, factor it as $t^\gamma a_0(1+r)$ with
$a_0\in\Hol(U)^\times$ and $v(r)>0$, then use the strong geometric series.
We write $\A^{>0}(U)$ for its positive-valuation part.

This ring is different from $\C\{w\}((t^\Gamma))$, whose coefficients are
individual convergent germs.  For example,
\[
 \sum_{n\ge1}\frac{t^{n\delta}}{1-nw}
\]
has convergent coefficients but no common ordinary disk about $0$.
None of the common-domain conclusions below is obtained by silently replacing
one ring by the other.

\subsection{Evaluation at surcomplex points}
\label{sec:evaluation}

Let $U^\#=\{a+\eta:a\in U,\ \eta\in\m\}$, where
$\m=\{x\in\K:v(x)>0\}$.  For
$\Phi=\sum_\gamma t^\gamma\phi_\gamma\in\A(U)$ define
\begin{equation}\label{eq:evaluation}
 \Phi^\#(a+\eta)
   =\sum_{\gamma\in\supp\Phi}\sum_{n\ge0}t^\gamma
       \frac{\phi_\gamma^{(n)}(a)}{n!}\eta^n.
\end{equation}
Lemma~\ref{lem:neumann}, with $T=\supp\Phi$ and
$S=\supp\eta$, proves summability.  Taylor's product identity and the same
finite-word argument show that evaluation is a ring homomorphism.
Thus the objects of this article do define functions at surcomplex points;
they are not merely ordinary functions with a suggestive parameter name.

The ordinary Taylor coefficients in~\eqref{eq:evaluation} can grow without
bound as $n$ or $\gamma$ varies.  This is harmless for strong summability,
which is controlled by exponents, not by numerical bounds on coefficients.
All derivatives in the article are derivatives in the spatial variable;
they are not a choice of derivation on the constant field of surreal
numbers.

\section{The strong logarithm of a near-identity substitution}
\label{sec:log}

We give the operator construction with its support certificate because it
is what permits the later theorem to work in arbitrary rank and on one
fixed ordinary domain.

\subsection{Taylor substitution and its support}

Let $\Delta\in\A^{>0}(U)$ and write $F=w+\Delta$.  Define
\begin{equation}\label{eq:substitution}
 C_F\Phi=\sum_{n\ge0}\frac{\Phi^{(n)}}{n!}\Delta^n,
 \qquad N=C_F-I.
\end{equation}
The derivatives in this expression are coefficientwise.  If
$S=\supp\Delta$ and $T=\supp\Phi$, every term has exponent in
$T+\Neu S$.  Lemma~\ref{lem:neumann} proves that the expression is a strong
sum.  Taylor's Leibniz formula proves that $C_F$ is a $\K$-linear ring
homomorphism.

Each term in $N^r\Phi$ contains at least $r$ factors drawn from $\Delta$
and its spatial derivatives.  Differentiation does not enlarge exponent
support.  Consequently, for any fixed exponent, only finitely many values
of $r$ and finitely many words of factors contribute.  This proves the
strong summability of all the series in the following lemma, even when
$\Gamma$ has no countable cofinal subset.

\begin{lemma}[A common-domain logarithmic vector field]
\label{lem:logvector}
The substitution $C_F$ is an automorphism of $\A(U)$.  Its strong logarithm
\begin{equation}\label{eq:operatorlog}
 D=\log C_F=\sum_{r\ge1}\frac{(-1)^{r+1}}r N^r
\end{equation}
is a $\K$-linear derivation, and there is a uniquely determined
$A\in\A^{>0}(U)$ such that
\begin{equation}\label{eq:vectorfield}
 D\Phi=A(w)\Phi'(w),\qquad A=D(w),\qquad \exp D=C_F.
\end{equation}
All coefficient functions of $A$ are holomorphic on the original $U$.
\end{lemma}

\begin{proof}
The inverse of $I+N$ is the strong series $\sum_{r\ge0}(-N)^r$.
The elementary one-variable identities for $\log(1+X)$ and $\exp X$ remain
valid on substituting $N$: each coefficient of each application involves
only finitely many words, so every rearrangement is finite at that
coefficient.

To prove the derivation assertion without assuming it, introduce an
ordinary scalar $s$ and the operator
\[
 P_s=\sum_{r\ge0}\binom{s}{r}N^r.
\]
For fixed inputs and a fixed Hahn exponent, both sides of
$P_s(\Phi\Psi)=P_s(\Phi)P_s(\Psi)$ are polynomials in $s$.  They agree at all
nonnegative integers, because then $P_s=C_F^s$ and $C_F$ is multiplicative.
They are therefore identical.  Differentiating at $s=0$ gives the Leibniz
rule, since
$\left.\frac{\dd}{\dd s}\binom{s}{r}\right|_{s=0}=(-1)^{r-1}/r$ for $r\ge1$.

It remains to justify the assertion that this derivation is a spatial
vector field.  For ordinary $\phi\in\Hol(U)$, each Hahn coefficient
$D_\gamma\phi$ is a $\C$-linear derivation from $\Hol(U)$ to itself.  At an
arbitrary $a\in U$, write
\[
 \phi(w)=\phi(a)+(w-a)\psi_a(w),\qquad \psi_a(a)=\phi'(a),
\]
where $\psi_a$ is holomorphic on all of $U$ by removable singularity.
Evaluation of the Leibniz rule at $a$ gives
\[
 (D_\gamma\phi)(a)=(D_\gamma w)(a)\phi'(a).
\]
This holds at every $a$, hence $D_\gamma\phi=(D_\gamma w)\phi'$.
The strong support certificate extends the identity coefficientwise to
arbitrary $\Phi\in\A(U)$.  It proves~\eqref{eq:vectorfield} and also proves
that no continuity assumption on an abstract derivation was smuggled into
the argument.
\end{proof}

\begin{remark}[No ordinary invariant-domain hypothesis]
The notation $C_F$ denotes the Taylor operation~\eqref{eq:substitution}.
It does not require that an ordinary numerical specialization of $F$ carry
$U$ into itself.  At a point $a+\eta\in U^\#$, its positive displacement
preserves standard part, so the evaluated substitution is legitimate.
The identity between evaluated substitution and composition follows by
finite-word Taylor regrouping.  It is not an appeal to an ordinary
self-map of $U$.
\end{remark}

\subsection{An exact identity at fixed points}

\begin{lemma}[Derivative of the logarithmic generator]
\label{lem:fixedderivative}
If an ordinary $a\in U$ satisfies $F(a)=a$, then
\begin{equation}\label{eq:fixedderivative}
 A(a)=0,\qquad A'(a)=\Log F'(a).
\end{equation}
\end{lemma}

\begin{proof}
Evaluation at $a$ annihilates $N\Phi$, so it annihilates $D(w)$.
If $\lambda_a=F'(a)$, the chain rule for Taylor substitution gives
\[
 (N\Phi)'(a)=(\lambda_a-1)\Phi'(a).
\]
Iteration gives $(N^r w)'(a)=(\lambda_a-1)^r$.  Apply the logarithm series.
All sums are strong because $\lambda_a-1$ has positive valuation, unless it
is zero, in which case the same identity is immediate.
\end{proof}

This exact identity, rather than a truncated modified differential
 equation, will determine the complete monodromy exponent.

\section{Support-controlled linearization on a common domain}
\label{sec:linearization}

The next theorem applies beyond polynomial Euler maps.  It records both the
leading ordinary dynamics and the precise support source for every
correction.

\begin{theorem}[Common-domain Hahn linearization]
\label{thm:common}
Let $U\subset\C$ be simply connected, $0\in U$, and let
$V\in\Hol(U)$ satisfy
\[
 V(0)=0,\qquad V'(0)=\mu\ne0,\qquad V(w)\ne0\quad(w\in U\setminus\{0\}).
\]
Let $E\in\A^{>0}(U)$ satisfy $E(0)=0$, and set
\begin{equation}\label{eq:generalF}
 F(w)=w+t^\delta\bigl(V(w)+E(w)\bigr),\quad
 \lambda=F'(0),\quad
 M=\Neu{\{\delta\}\cup\supp E}.
\end{equation}
There is a unique $H\in\A(U)$ with
\begin{equation}\label{eq:generalH}
 H\circ F=\lambda H,\qquad H(0)=0,
 \qquad H'(0)=1.
\end{equation}
It has support contained in $M$, and every coefficient is holomorphic on
$U$.  Its coefficient at exponent zero is
\begin{equation}\label{eq:leadingH}
 h_0(w)=w\exp\!\left(\int_0^w
        \left(\frac{\mu}{V(\zeta)}-\frac1\zeta\right)\dd\zeta\right).
\end{equation}
All constructions commute with enlargement of the set-sized value group.
\end{theorem}

\begin{proof}
Lemma~\ref{lem:logvector} gives $\log C_F=A\partial_w$.
Every word in $\Delta=t^\delta(V+E)$ has an exponent which, after removing
one initial $\delta$, belongs to $M$.  Therefore
\[
 A=t^\delta W,\qquad W=V+W_+,\qquad
 \supp W\subset M,
\]
where $W_+$ has positive valuation and all coefficients are holomorphic
on $U$.  Set
\[
 c=\Log\lambda,\qquad \alpha=t^{-\delta}c.
\]
Lemma~\ref{lem:fixedderivative} gives $W(0)=0$ and $W'(0)=\alpha$.
The leading coefficient of $\alpha$ is $\mu$.

Coefficientwise division by $w$ is possible since $W(0)=0$.
The leading coefficient of $W/w$ is $V/w$, which is nowhere zero on $U$,
including at $0$.  Thus $W/w$ is invertible in $\A(U)$.  Moreover,
\begin{equation}\label{eq:regularR}
 R(w):=\frac{\alpha}{W(w)}-\frac1w
   =\frac{\alpha/(W(w)/w)-1}{w}
\end{equation}
belongs to $\A(U)$: the numerator vanishes coefficientwise at $0$ because
$(W/w)(0)=\alpha$.  Its support is contained in $M$.

Every coefficient of $R$ has a unique primitive vanishing at $0$, since $U$
is simply connected.  Denote the resulting Hahn series by
$B(w)=\int_0^w R(\zeta)\dd\zeta$.  Write $B=B_0+B_+$, where $B_0$ is the
ordinary coefficient at zero.  Define
\begin{equation}\label{eq:Hconstruction}
 H(w)=w e^{B_0(w)}\sum_{n\ge0}\frac{B_+(w)^n}{n!}.
\end{equation}
Lemma~\ref{lem:neumann} certifies the exponential and preserves the support
bound $M$.  This construction uses the same $U$ at every stage.  It gives
$H(0)=0$ and $H'(0)=1$, and differentiation gives
\[
 \frac{H'}H=\frac\alpha W
\]
away from $0$, with the corresponding cross-multiplied identity everywhere.
Thus $D H=cH$.  As $D$ fixes Hahn constants, exponentiation yields
$C_FH=\ExpH(c)H=\lambda H$.
Formula~\eqref{eq:leadingH} is the exponent-zero part of the construction.

For uniqueness, let $\widetilde H$ satisfy~\eqref{eq:generalH}.
Then $N^r\widetilde H=(\lambda-1)^r\widetilde H$, whence
$D\widetilde H=c\widetilde H$.  The quotient
$\widetilde H/H=(\widetilde H/w)/(H/w)$ is in $\A(U)$ because the leading
coefficient of $H/w$ is the nowhere-zero function $e^{B_0}$.
The differential equations imply that its derivative is zero.  For example,
multiplication by the nonzero series $W$ gives zero derivative first, and
$\A(U)$ is an integral domain since $U$ is connected.  Every coefficient of
the quotient is consequently constant.  Its value at $0$ is $1$, by the
derivative normalization, so $\widetilde H=H$.

Finally, all outputs are formed from input supports by finite words,
ordinary coefficientwise differentiation, and uniquely normalized
integration.  None of these depends on the ambient group beyond containing
the indicated supports.  This proves compatibility under enlargement.
\end{proof}

\begin{remark}[What kind of embedding was proved]
The proof constructs a strong logarithm and solves a coefficientwise
ordinary differential equation.  It does not assert convergence of the
series in a nonzero ordinary complex value of $t^\delta$, nor analytic
embedding of every numerical near-identity map in an autonomous flow.
Those are different questions, and the distinction is essential
\cite{GV}.  The theorem also does not assert that the global map $H^\#$ is
injective on all of $U^\#$; injectivity on an appropriate infinitesimal disk
will be proved separately for the resonant model.
\end{remark}

\subsection{A universality consequence for higher analytic terms}

\begin{corollary}[Universal leading profile at a resonant scale]
\label{cor:universal}
Let $q\ge1$, $a\in\C^*$, and let
$\phi(z)=a z^{q+1}+\sum_{n>q+1}a_nz^n$ be an ordinary convergent germ.
Choose $c\in\C^*$ with $a c^q=1$, set $s=c\,t^{\delta/q}$, and consider
$f(z)=(1+\eps)z+\phi^\#(z)$ at infinitesimal $z$.
After $z=sw$, the normalized Hahn coordinate on any simply connected
ordinary $U$ containing $0$ but no root of $1+w^q$ has leading coefficient
\begin{equation}\label{eq:universalprofile}
 w(1+w^q)^{-1/q},
\end{equation}
with the branch normalized at $0$.
\end{corollary}

\begin{proof}
Taylor evaluation of the convergent germ at $sw$ is strongly meaningful
for every finite $w$.  Direct rescaling gives
\[
 s^{-1}f(sw)=w+\eps\left(w+w^{q+1}
   +\sum_{n>q+1} a_nc^{n-1}
       t^{\delta(n-q-1)/q}w^n\right).
\]
The final sum has positive Hahn support and polynomial coefficient
functions on any such $U$.  Apply Theorem~\ref{thm:common} with
$V=w(1+w^q)$; integration of $1/V-1/w$ gives
$-q^{-1}\log(1+w^q)$.
\end{proof}

The corollary asserts a leading profile and a common-domain construction.
It does not extend the exact monodromy formula or the sharp disk theorem
below to arbitrary higher-order perturbations without additional work.

\section{Polynomial Euler maps and their exact logarithmic residues}
\label{sec:polynomial}

We now specialize to
\begin{equation}\label{eq:Euler}
 G_\eps(w)=w+\eps V(w),\qquad V\in\C[w],
 \qquad V(0)=0,\quad V'(0)=\mu\ne0.
\end{equation}
Theorem~\ref{thm:common}, with $E=0$, constructs $H_\eps$ in
$\Hol(U)[[\eps]]$.  The polynomial hypothesis provides additional structure
which is not part of the general theorem.

\subsection{The generator has no new zero divisor}

\begin{proposition}[Polynomial factorization of the generator]
\label{prop:factorV}
The logarithmic generator of~\eqref{eq:Euler} has the form
\begin{equation}\label{eq:Wfactor}
 \log C_{G_\eps}=\eps W_\eps(w)\partial_w,
 \qquad W_\eps(w)=V(w)B_\eps(w),
 \qquad B_\eps\in1+\eps\C[w][[\eps]].
\end{equation}
Set
\begin{equation}\label{eq:alphaOmega}
 \alpha_\eps=\frac{\Log(1+\mu\eps)}\eps,
 \qquad \Omega_\eps(w)=\frac{\alpha_\eps}{W_\eps(w)}.
\end{equation}
Every coefficient of $\Omega_\eps$ is a rational function of $w$ with poles
only at zeros of $V$.  The order of a pole is at most its multiplicity in
$V$.
\end{proposition}

\begin{proof}
For a polynomial $p$,
\[
 (C_{G_\eps}-I)p
   =\sum_{k\ge1}\frac{\eps^k V^k}{k!}p^{(k)}.
\]
Every coefficient is polynomial and divisible by $V$.  The same is true of
each $N^r w$, and the coefficient at any fixed power of $\eps$ involves
only finitely many $r$.  Thus the logarithm has polynomial coefficients,
all divisible by $V$.  Its first term is $\eps V$, proving
\eqref{eq:Wfactor}.  The inverse of $B_\eps$ has polynomial coefficients by
the formal geometric series.  Consequently $\Omega_\eps$ is a formal series
of polynomials divided by the single polynomial $V$.
\end{proof}

The divisor assertion is stronger than merely knowing that each coefficient
is rational.  Higher terms of the modified vector field cannot introduce
new spatial pole locations into the logarithmic derivative of the
coordinate.

\begin{lemma}[Exact residue at a simple fixed point]
\label{lem:residue}
At a simple zero $a$ of $V$, with $\nu_a=V'(a)$,
\begin{equation}\label{eq:exactresidue}
 \Res_{w=a}\bigl(\Omega_\eps(w)\dd w\bigr)
 =\rho_a(\eps)
 =\frac{\Log(1+\mu\eps)}{\Log(1+\nu_a\eps)}.
\end{equation}
In particular $\rho_0=1$.
\end{lemma}

\begin{proof}
Lemma~\ref{lem:fixedderivative}, applied at $a$, gives
\[
 \eps W_\eps'(a)=\Log G_\eps'(a)=\Log(1+\nu_a\eps).
\]
Since $a$ is a simple zero of $V$, the factorization
$W_\eps=V B_\eps$ makes it a simple zero over $\C[[\eps]]$ as well.
The residue of $\alpha_\eps/W_\eps$ is
$\alpha_\eps/W_\eps'(a)$, which is~\eqref{eq:exactresidue}.
\end{proof}

These are coefficientwise ordinary residues of a rational differential.
They are not local residues at displaced surcomplex poles.  The fixed
points in this polynomial family are exactly the ordinary zeros of $V$,
so there is no displacement ambiguity in the present calculation.

\subsection{Modified-vector-field coefficients and the first correction}

Finite expansion of the operator logarithm gives
\begin{equation}\label{eq:modifiedvector}
 W_\eps
 =V-\frac\eps2 VV'
 +\eps^2\left(\frac1{12}V^2V''+\frac13 V(V')^2\right)
 +O(\eps^3).
\end{equation}
Here and below $O(\eps^k)$ means divisibility by $\eps^k$ in the indicated
formal series ring; it is not an analytic estimate uniform in a numerical
parameter.

For completeness, $Nw=\eps V$ and
\[
 N^2w=\eps^2 VV'+\frac{\eps^3}{2}V^2V''+O(\eps^4),
 \qquad
 N^3w=\eps^3\bigl(V(V')^2+V^2V''\bigr)+O(\eps^4).
\]
Substitution in $N-N^2/2+N^3/3$ proves~\eqref{eq:modifiedvector}.
No assertion about all coefficients is being inferred from this truncation;
the exact identity~\eqref{eq:exactresidue} has already been proved separately.

Let $h_0$ be the leading coordinate~\eqref{eq:leadingH}.  The quotient
$V/(\mu h_0)$ is holomorphic and nowhere zero on $U$, and equals $1$ at $0$
after removing the common zero.  Its logarithm is therefore uniquely fixed
by requiring value $0$ at $0$.

\begin{proposition}[The first correction is an explicit logarithm]
\label{prop:firstcorrection}
On the common domain $U$,
\begin{equation}\label{eq:firstcorrection}
 H_\eps
  =h_0\left[1+\frac{\mu\eps}{2}
       \log\!\left(\frac{V}{\mu h_0}\right)+O(\eps^2)\right].
\end{equation}
\end{proposition}

\begin{proof}
Since $\alpha_\eps=\mu-\mu^2\eps/2+O(\eps^2)$,
formula~\eqref{eq:modifiedvector} gives
\[
 \Omega_\eps=\frac\mu V
       +\frac{\mu\eps}{2}\frac{V'-\mu}{V}+O(\eps^2).
\]
The coefficient at $\eps$ is holomorphic at $0$ and equals
$\frac\mu2\bigl(\log(V/(\mu h_0))\bigr)'$, because
$h_0'/h_0=\mu/V$.  Integrate with the normalization at $0$ and exponentiate.
\end{proof}

\subsection{A global factorization when all zeros are simple}

Suppose $V$ has degree $d\ge2$ and distinct zeros $0,a_1,\ldots,a_{d-1}$.
Partial fractions, taken separately at every $\eps$ coefficient, give
\begin{equation}\label{eq:partialfractions}
 \Omega_\eps(w)=Q_\eps(w)+\frac1w
      +\sum_{j=1}^{d-1}\frac{\rho_{a_j}(\eps)}{w-a_j},
 \qquad Q_\eps\in\C[w][[\eps]].
\end{equation}
The coefficients at $\eps^0$ and $\eps^1$ have numerator degree smaller
than $d$, by the preceding formulas.  Hence $Q_\eps\in\eps^2\C[w][[\eps]]$.
Put $P_\eps(w)=\int_0^w Q_\eps(\zeta)\dd\zeta$.

\begin{corollary}[Separation of the branching factors]
\label{cor:factorization}
On branches normalized near $0$,
\begin{equation}\label{eq:fullfactorization}
 H_\eps(w)=w\ExpH(P_\eps(w))
      \prod_{j=1}^{d-1}(1-w/a_j)^{\rho_{a_j}(\eps)},
 \qquad P_\eps\in\eps^2\C[w][[\eps]],\quad P_\eps(0)=0.
\end{equation}
The factor $\ExpH(P_\eps)$ and its inverse have polynomial coefficients
and belong to $\B$.  In particular all nontrivial coefficientwise
monodromy is carried by the displayed power factors.
\end{corollary}

\begin{proof}
For $\rho=\rho_0+\rho_+$, define the power on a chosen ordinary logarithm
branch by
\[
 (1-w/a)^\rho
  =\exp\bigl(\rho_0\log(1-w/a)\bigr)
     \ExpH\bigl(\rho_+\log(1-w/a)\bigr).
\]
The logarithmic derivative of the right side of
\eqref{eq:fullfactorization} is~\eqref{eq:partialfractions}; its value and
derivative at $0$ have the required normalization.  Uniqueness proves the
identity.  At a fixed power of $\eps$, the exponential of $P_\eps$ uses
only finitely many products of polynomials, proving the final assertion.
\end{proof}

\section{Infinitesimal monodromy and the finite-cover obstruction}
\label{sec:monodromy}

\subsection{Analytic continuation in the correct ambient field}

Let $X=\C\setminus Z(V)$ and let $p:\Ucov\to X$ be its ordinary universal
cover.  Choose a small nonzero base point near $0$ to transport the
normalized germ.  Every coefficient of $\Omega_\eps$ is holomorphic on $X$.
Its coefficientwise primitive on $\Ucov$, followed by exponentiation, gives
a continuation of every coefficient of $H_\eps$.
The leading coefficient has no zeros on $\Ucov$.

The common ambient field for algebraicity is
\begin{equation}\label{eq:ambientfield}
 \mathscr L=\Mer(\Ucov)((t^\Gamma)),
 \qquad \B=\C(w)((t^\Gamma))\hookrightarrow\mathscr L,
\end{equation}
where rational functions are pulled back along $p$.  It is important to use
\emph{meromorphic} coefficients in this ambient field.  An arbitrary element
of $\B$ can have rational coefficient functions whose additional poles vary
with the exponent; no one common pole-free subdomain for all elements of
$\B$ is assumed.  Each pulled-back coefficient is nevertheless
meromorphic on $\Ucov$, which is exactly what~\eqref{eq:ambientfield}
requires.

A deck transformation acts on each coefficient by analytic continuation
and leaves Hahn exponents unchanged.  It is therefore a strong field
automorphism of $\mathscr L$, and fixes every element of $\B$.
This is a statement about an ordinary covering space and a field of
coefficient functions.  No covering-space topology on $\No[i]$ is being
introduced.

\subsection{The torsion test for a finite Hahn exponential}

\begin{lemma}[An infinitesimal phase has no torsion]
\label{lem:torsion}
For a finite Hahn constant $c=c_0+c_+$, with $c_0\in\C$ and $v(c_+)>0$,
\[
 \ExpH(c)=1\quad\Longleftrightarrow\quad c_+=0
       \text{ and }c_0\in2\pi i\Z.
\]
Consequently, for $\rho\in\C+\m$, the element $\ExpH(2\pi i\rho)$ has
finite multiplicative order if and only if $\rho$ is an ordinary rational
number.
\end{lemma}

\begin{proof}
The coefficient at exponent zero in $\ExpH(c)$ is $e^{c_0}$, so equality
to $1$ first requires $c_0\in2\pi i\Z$.  If $c_+\ne0$, then
$\ExpH(c_+)-1=c_++c_+^2/2+\cdots$ has the same leading term as $c_+$:
all higher powers have strictly larger valuation.  It cannot vanish.
For the last assertion, the exponential product law gives
$\ExpH(2\pi i\rho)^m=\ExpH(2\pi i m\rho)$, and the first assertion applies.
\end{proof}

\begin{theorem}[Exact monodromy and obstruction to finite ramification]
\label{thm:monodromy}
At every simple zero $a$ of $V$, continuation of $H_\eps$ around a positive
simple loop about $a$ is multiplication by $M_a$ in
\eqref{eq:monodromyintro}.  The multiplier has finite order exactly when
$V'(a)=\mu$, and then it is $1$.  When $V'(a)\ne\mu$, no finite ordinary
branched cover at $a$ makes this continued coordinate single-valued.
\end{theorem}

\begin{proof}
Near $a$, the rational-coefficient differential decomposes as
\[
 \Omega_\eps(w)\dd w
     =\rho_a\frac{\dd w}{w-a}+\text{a holomorphic-coefficient differential}.
\]
A positive circuit changes a coefficientwise primitive by $2\pi i\rho_a$.
Exponentiation therefore multiplies $H_\eps$ by $M_a$.
Write $\nu=V'(a)$.  Formal division of the two logarithms gives
\begin{equation}\label{eq:rhofirst}
 \rho_a=\frac\mu\nu+
          \frac{\mu(\nu-\mu)}{2\nu}\eps+O(\eps^2).
\end{equation}
If $\nu\ne\mu$, the coefficient of $\eps$ is nonzero, and
Lemma~\ref{lem:torsion} makes $M_a$ infinite order.  If $\nu=\mu$, the ratio
of logarithms is exactly $1$.

A finite branched cover has a positive integer ramification index $m$ at a
point above $a$.  In suitable local coordinates, a circuit upstairs
projects to $m$ circuits downstairs.  Its monodromy multiplier is $M_a^m$,
which is not $1$ in the infinite-order case.  Thus no such cover removes
the branching.  This applies in particular to the covers used to make an
ordinary algebraic function single-valued.
\end{proof}

The obstruction is already first order in $\eps$ whenever $\nu\ne\mu$.
It is not necessary to inspect infinitely many unknown correction terms to
decide whether finite ramification will work.

\subsection{Why some simple root must obstruct}

\begin{lemma}[A nonlinear polynomial cannot have the same slope at all its
simple zeros]\label{lem:slopes}
If a polynomial of degree $d\ge2$ has $d$ distinct zeros, then
\begin{equation}\label{eq:residuesum}
 \sum_{V(a)=0}\frac1{V'(a)}=0.
\end{equation}
In particular, if $V'(0)=\mu\ne0$, not all slopes at its zeros can equal
$\mu$.
\end{lemma}

\begin{proof}
The partial-fraction expansion is
$1/V(w)=\sum_a(V'(a)(w-a))^{-1}$.  The coefficient of $1/w$ at infinity
on the left is zero because $d\ge2$; on the right it is the indicated sum.
If every derivative were $\mu$, the sum would be $d/\mu\ne0$.
\end{proof}

Thus Theorem~\ref{thm:monodromy} supplies at least one infinite-order
multiplier in every nonlinear simple-root case, even when some individual
roots have trivial monodromy.

\section{The Hahn-algebraicity dichotomy}
\label{sec:algebraicity}

We now prove the algebraicity assertion of Theorem~\ref{thm:main}, including
the multiple-root case not covered by simple-pole monodromy.

\subsection{Two elementary algebraicity tests}

\begin{lemma}[Reduction of an algebraic Hahn element]
\label{lem:reduction}
Suppose $Y\in\mathscr L$ has valuation zero and leading coefficient $y_0$.
If $Y$ is algebraic over $\B$, then $y_0$ is algebraic over $\C(w)$.
\end{lemma}

\begin{proof}
Take a nonzero polynomial relation $\sum_{j=0}^m b_jY^j=0$ with
$b_j\in\B$.  Let $\gamma$ be the least valuation among the nonzero $b_j$.
Multiply the relation by $t^{-\gamma}$.  Now all coefficients have
nonnegative valuation and at least one has nonzero coefficient at zero.
Taking coefficient at zero yields a nonzero polynomial over $\C(w)$
annihilating $y_0$.  Multiplication reduces correctly because all the
series in this normalized relation have nonnegative valuation.
\end{proof}

\begin{lemma}[A higher-order logarithmic pole forces transcendence]
\label{lem:essential}
Let $r(w)$ be a rational function with a pole of order at least two at
$a$.  Any nonzero local solution of $y'/y=r(w)$ is transcendental over
$\C(w)$.
\end{lemma}

\begin{proof}
On a small punctured sector, integration gives
\[
 y(w)=u(w)(w-a)^\beta
       \exp\!\left(\sum_{j=1}^{k-1}c_j(w-a)^{-j}\right),
 \qquad c_{k-1}\ne0,
\]
where $u$ is holomorphic and nonzero and $k\ge2$ is the pole order of $r$.
Along a suitable ray ending at $a$, the real part of the leading expression
$c_{k-1}(w-a)^{-(k-1)}$ is positive.  It dominates the lower terms, so
$|y(w)|$ grows faster than every power of $|w-a|^{-1}$ along that ray.

An algebraic function cannot do this.  Indeed, after dividing a polynomial
relation by its leading coefficient, its coefficients are meromorphic and
hence bounded by some fixed power of $|w-a|^{-1}$ on a sufficiently small
punctured disk.  The elementary bound
$|y|\le1+\max_{j<m}|b_j|$ for a root of
$y^m+\sum_{j<m}b_j y^j$ gives a polynomial growth bound for every root.
This contradicts the preceding ray estimate.
\end{proof}

\subsection{Proof of the principal dichotomy}

\begin{proof}[Proof of Theorem~\ref{thm:main}]
Existence, uniqueness, and the common-domain assertion are
Theorem~\ref{thm:common} with $E=0$.  Continuation and the monodromy formula
were established in Sections~\ref{sec:polynomial}--\ref{sec:monodromy}.
If $V$ is linear, it is $\mu w$, and $H_\eps=w$ solves the equation; it
belongs to $\B$.

Assume henceforth that $\deg V\ge2$.  If $V$ has a multiple zero, then
$\mu/V$ has a pole of order at least two there.  The leading coordinate
satisfies $h_0'/h_0=\mu/V$, so Lemma~\ref{lem:essential} makes $h_0$
transcendental over $\C(w)$.  Since $H_\eps$ has valuation zero as an
element of $\mathscr L$, Lemma~\ref{lem:reduction} excludes algebraicity
over $\B$.

If every zero is simple, Lemma~\ref{lem:slopes} supplies a zero $a$ with
$V'(a)\ne\mu$.  Let $\tau_a$ be the corresponding deck action on
$\mathscr L$.  It fixes $\B$ and satisfies
\[
 \tau_a^n(H_\eps)=M_a^n H_\eps\qquad(n\ge0).
\]
Theorem~\ref{thm:monodromy} shows that these elements are all distinct;
$H_\eps\ne0$ and $M_a$ has infinite order.  If $H_\eps$ were algebraic over
$\B$, all these elements would be roots of one nonzero polynomial over
$\B$, which is impossible in a field.  This proves the reverse implication
in~\eqref{eq:dichotomyintro} and completes the theorem.
\end{proof}

The argument works over the larger field $\B$ precisely because ordinary
deck transformations fix \emph{every rational coefficient separately}.
No degree bound on those rational functions, no finite support assumption,
and no common bound on their pole sets is needed.

\subsection{Exactly when the leading answer is misleading}

\begin{proposition}[Classification of algebraic leading coordinates]
\label{prop:leadingalgebraic}
For a nonlinear polynomial $V$ as above, $h_0$ is algebraic over $\C(w)$
if and only if all zeros of $V$ are simple and
\begin{equation}\label{eq:rationalresidues}
 \frac\mu{V'(a)}\in\Q\qquad\text{for every zero }a\text{ of }V.
\end{equation}
Whenever these conditions hold, the first correction $h_1$ in
$H_\eps=h_0+\eps h_1+\cdots$ is already transcendental over $\C(w)$.
\end{proposition}

\begin{proof}
Multiple zeros are excluded by Lemma~\ref{lem:essential}.  With all zeros
simple, the leading part of~\eqref{eq:fullfactorization} is
\begin{equation}\label{eq:h0product}
 h_0(w)=w\prod_{a\ne0}(1-w/a)^{\mu/V'(a)}.
\end{equation}
If all exponents are rational, taking a common denominator gives a
polynomial relation, so $h_0$ is algebraic.  Conversely, continuation of an
algebraic function gives a finite orbit.  The local multipliers of
\eqref{eq:h0product} must therefore have finite order.  For an ordinary
complex exponent this is equivalent to its being rational, proving the
classification.

For the last assertion, choose $a$ with $\nu=V'(a)\ne\mu$ by
Lemma~\ref{lem:slopes}.  Put $r=\mu/\nu\in\Q$ and choose a positive integer
$m$ with $mr\in\Z$.  After $m$ circuits about $a$, $h_0$ returns to itself.
By Proposition~\ref{prop:firstcorrection}, the same continuation changes
$h_1$ by
\[
 2\pi i m\,\frac\mu2(1-r)h_0,
\]
which is nonzero.  Repeated blocks of $m$ circuits therefore produce
infinitely many distinct branches of $h_1$.  An algebraic function has
only finitely many branches, so $h_1$ is transcendental.
\end{proof}

\begin{example}[A rational leading coordinate and a logarithmic correction]
For $V=w+w^2$, $h_0=w/(1+w)$ and
\[
 h_1=\frac{w}{1+w}\log(1+w).
\]
At the fixed point $-1$, the exact exponent is
$\Log(1+\eps)/\Log(1-\eps)=-1+\eps+O(\eps^2)$.
Thus the leading multiplier is $1$, while the full multiplier begins
$1+2\pi i\eps+O(\eps^2)$.  The obstruction is present even without any
branching in the leading function.
\end{example}

\begin{example}[A nonlinear map with one trivial nonzero-root multiplier]
Let $V=w(1-w)(1-2w)$.  Then $\mu=1$, $V'(1)=1$, and
$V'(1/2)=-1/2$.  The leading coordinate and first correction are
\[
 h_0=\frac{w(1-w)}{(1-2w)^2},\qquad
 h_1=\frac32 h_0\log(1-2w).
\]
The exact multiplier at $1$ is $1$, whereas at $1/2$ it is
$1+3\pi i\eps+O(\eps^2)$ and has infinite order.
Thus the theorem does not claim branching at every nonzero fixed point.
The slope identity guarantees at least one obstruction, which is sufficient
for transcendence over the whole Hahn rational-function base.
\end{example}

\begin{example}[Why multiple zeros require a separate argument]
For $V=w(1+w)^2$, $\mu=1$ and
\[
 h_0(w)=\frac{w}{1+w}\exp\!\left(\frac1{1+w}-1\right).
\]
The pole of $\mu/V$ at $-1$ has order two.  The leading solution has an
essential singularity and is already transcendental.  A proof discussing
only simple-pole residues would miss this case, even though the full
dichotomy remains true.
\end{example}

\section{The resonant model: exact scale, isometry, and continuation}
\label{sec:model}

Fix an ordinary integer $q\ge1$ and consider
\[
 f(z)=(1+\eps)z+z^{q+1},\qquad
 \eps=t^\delta,\quad s=t^{\delta/q}.
\]
The divisibility hypothesis on $\Gamma$ ensures that $s\in\K$.
Under $z=sw$, the map becomes
\begin{equation}\label{eq:gmodel}
 g(w)=s^{-1}f(sw)=(1+\eps)w+\eps w^{q+1}
              =w+\eps V_q(w),\quad V_q(w)=w(1+w^q).
\end{equation}
Let $\Psi_\eps$ denote its normalized common-domain coordinate.
The disk $U=\{|w|<1\}$ is one admissible domain; larger simply connected
slit domains avoiding the roots of $1+w^q$ are also possible.

\subsection{Exact coefficients and their valuations}

For a $q$th root of unity $\zeta$, the map obeys
$g(\zeta w)=\zeta g(w)$.  Uniqueness of the normalized coordinate implies
$\Psi_\eps(\zeta w)=\zeta\Psi_\eps(w)$.  Therefore its Taylor expansion has
the form
\begin{equation}\label{eq:bseries}
 \Psi_\eps(w)=\sum_{m\ge0}b_m(\eps)w^{qm+1},
 \qquad b_0=1,\quad b_m\in\C[[\eps]].
\end{equation}
This initially means the Taylor series of each ordinary coefficient
function; its evaluation as a Hahn family at a specific $w$ will be a
separate question.

The leading solution is
\begin{equation}\label{eq:leadingmodel}
 \Psi_0(w)=w(1+w^q)^{-1/q},
 \qquad
 b_m(0)=(-1)^m\frac{(1/q)_m}{m!}\ne0,
\end{equation}
where $(a)_m=a(a+1)\cdots(a+m-1)$ and $(a)_0=1$.

\begin{proposition}[A finite recurrence]
\label{prop:recurrence}
Write $\lambda=1+\eps$.  For $m\ge1$,
\begin{equation}\label{eq:brecurrence}
 b_m=-\frac{
  \displaystyle\sum_{\substack{0\le j<m\\m-j\le qj+1}}
     b_j\binom{qj+1}{m-j}
       \lambda^{qj+1-(m-j)}\eps^{m-j}}
 {\lambda^{qm+1}-\lambda}.
\end{equation}
This recurrence uniquely constructs $b_m\in\C[[\eps]]$ and has leading
recurrence
\begin{equation}\label{eq:leadingrecurrence}
 qm\,b_m(0)=-\bigl(q(m-1)+1\bigr)b_{m-1}(0).
\end{equation}
\end{proposition}

\begin{proof}
In $g(w)^{qj+1}=w^{qj+1}(\lambda+\eps w^q)^{qj+1}$, the coefficient of
$w^{qm+1}$ is the binomial expression in~\eqref{eq:brecurrence}.
The term $j=m$ contributes $\lambda^{qm+1}b_m$, and the right side of
$\Psi_\eps\circ g=\lambda\Psi_\eps$ contributes $\lambda b_m$.
This proves the recurrence.  Its denominator is
$qm\eps+O(\eps^2)$, while every term in the numerator is divisible by
$\eps$.  Induction therefore gives $b_m\in\C[[\eps]]$.
Only $j=m-1$ survives after division by $\eps$ and reduction modulo $\eps$,
which gives~\eqref{eq:leadingrecurrence} and also checks
\eqref{eq:leadingmodel} directly.
\end{proof}

The normalized formal coordinate for the original map is
\begin{equation}\label{eq:originalh}
 h(z)=s\Psi_\eps(z/s)
      =\sum_{m\ge0}\eps^{-m}b_m(\eps)z^{qm+1}.
\end{equation}
This is a formal power series in $z$ over $\K$, not a Hahn series whose
coefficient support has already been asserted summable at every $z$.
Its nonzero coefficients have the exact valuations
\begin{equation}\label{eq:coefvaluation}
 v\bigl(\eps^{-m}b_m(\eps)\bigr)=-m\delta.
\end{equation}
The ordinary formal Schr\"oder equation also gives uniqueness directly,
since $\lambda^n-\lambda\ne0$ for every integer $n\ge2$.

\subsection{The exact strong summability disk}

\begin{theorem}[Sharp domain and valuative isometry]
\label{thm:sharp}
The family of Taylor terms in~\eqref{eq:originalh} is strongly summable at
$z\in\K$ if and only if $z=0$ or
\begin{equation}\label{eq:sharpdomain}
 v(z)>\delta/q.
\end{equation}
On the disk $D=\{z:v(z)>\delta/q\}$, the resulting function is a bijection
$h:D\to D$, satisfies $h\circ f=\lambda h$, and obeys
\begin{equation}\label{eq:isometry}
 v(h(x)-h(y))=v(x-y)\qquad(x,y\in D).
\end{equation}
No strictly larger centered valuative ball containing $D$ admits an
injective conjugacy to multiplication by $\lambda$ which sends $0$ to $0$.
\end{theorem}

\begin{proof}
Let $z\ne0$, $\beta=v(z)$.  The valuation of its $m$th nonzero Taylor term
is exactly
\begin{equation}\label{eq:termvaluation}
 v\bigl(\eps^{-m}b_m z^{qm+1}\bigr)
       =\beta+m(q\beta-\delta).
\end{equation}
If $q\beta=\delta$, every term has a nonzero coefficient at the same
exponent $\beta$.  Local finiteness fails.  If $q\beta<\delta$, the displayed
leading exponents form a strictly descending sequence in the union of
supports.  Well-ordering fails.  This proves necessity without any
convergence terminology.

If $q\beta>\delta$, put $w=z/s$, so $v(w)>0$.  In~\eqref{eq:bseries},
all coefficients $b_m$ have nonnegative $\eps$-support.  The support of its
Taylor family is governed by the positive well-ordered set
$\{\delta\}\cup\supp w$.  A term of spatial degree $qm+1$ contributes a
word of at least that length from $\supp w$, together with a nonnegative
number of $\delta$'s.  Lemma~\ref{lem:neumann} makes the full double family
strongly summable.  Multiplication by $s$ preserves this property, proving
sufficiency and making substitution in the formal conjugacy identity legal.

For invertibility, formal reversion of
$\Psi_\eps(w)=w+O(w^{q+1})$ gives an inverse
$\Theta_\eps\in\C[[\eps]][[w]]$ with linear coefficient $1$ and no
constant term.  This requires no division except by the unit linear
coefficient; every coefficient is obtained from finitely many preceding
coefficients.  The same positive-word argument evaluates $\Theta_\eps$
at every $w\in\m$.  The formal identities
$\Theta_\eps\circ\Psi_\eps=w$ and
$\Psi_\eps\circ\Theta_\eps=w$ become strong identities there.
Both maps preserve $\m$, since their difference from the identity has
strictly greater valuation at a nonzero infinitesimal input.  Scaling by
$s$ proves $h:D\to D$ bijective.

For the isometry, put $u=x/s$ and $v_1=y/s$, both of positive valuation.
Factoring the difference of powers gives
\[
 \Psi_\eps(u)-\Psi_\eps(v_1)
  =(u-v_1)\left[1+\sum_{m\ge1}b_m(\eps)
       \sum_{j=0}^{qm}u^{qm-j}v_1^j\right].
\]
The bracket is a strong sum equal to $1$ plus a positive-valuation element;
each nonconstant summand has valuation at least
$q\min\{v(u),v(v_1)\}>0$ when both inputs are nonzero.  A zero input is
handled by the identical argument with its vanishing terms omitted.
Thus the bracket has valuation zero, proving~\eqref{eq:isometry} after
rescaling.  The case $x=y$ is immediate with $v(0)=+\infty$.

Finally, every nonzero solution of $z^q=-\eps$ is a fixed point of $f$ and
has valuation $\delta/q$.  Such roots belong to $\K$, since their constant
factors are complex $q$th roots of $-1$.  Any strictly larger centered
valuative ball containing $D$ contains all elements at valuation
$\delta/q$, hence at least one such fixed point in addition to $0$.
An injective conjugacy $k$ would satisfy
$k(z)=\lambda k(z)$ at every fixed point.  As $\lambda\ne1$, it would map
both fixed points to $0$, a contradiction.
\end{proof}

The final assertion concerns centered balls defined by a valuation
threshold, including open or closed boundary conventions.  It does not
exclude continuation to a non-ball domain which avoids the obstructing
fixed points, and it does not assert that every coefficientwise continuation
is injective.  These distinctions matter at the critical scale.

\subsection{Continuation at a scale where the Taylor family fails}

Let $w_0\in U$ be an ordinary nonzero point.  At $z=sw_0$ we have
$v(z)=\delta/q$, so the Taylor family~\eqref{eq:originalh} is not strongly
summable.  On the other hand,
\begin{equation}\label{eq:resummedvalue}
 s\Psi_\eps(w_0)=s\sum_{n\ge0}\eps^n\Psi_n(w_0)
\end{equation}
is a valid Hahn value.  More generally it is defined at $z=s(w_0+\eta)$
for positive $\eta$, using~\eqref{eq:evaluation}.

There is no contradiction.  In~\eqref{eq:resummedvalue}, each coefficient
function $\Psi_n$ is first an ordinary holomorphic function, obtained by
ordinary integration and continuation.  Summation in the Hahn exponent
comes afterward.  In~\eqref{eq:originalh}, summation in spatial Taylor degree
is instead required to be strongly summable as a family of Hahn numbers,
and it fails precisely on this boundary.
The boundary value is a \emph{coefficientwise analytic continuation}, not an
illegal regrouping of a strong sum.  On the inner disk both procedures are
admissible and agree by finite-word Taylor identities.

This distinction gives the exact setting for the monodromy phenomenon:
ordinary continuation moves among branches of coefficient functions at the
rescaled boundary, even though the original Taylor representation stops
there.

\section{Explicit hidden monodromy at the critical scale}
\label{sec:explicit}

For $V_q=w(1+w^q)$, every nonzero zero $a$ satisfies $a^q=-1$ and
$V_q'(a)=-q$.  Hence all these zeros have the same exact exponent
\begin{equation}\label{eq:rhoq}
 \rho_q(\eps)=\frac{\Log(1+\eps)}{\Log(1-q\eps)}
  =-\frac1q+A_q\eps+B_q\eps^2+O(\eps^3),
\end{equation}
where
\begin{equation}\label{eq:ABC}
 A_q=\frac{q+1}{2q},\qquad
 B_q=\frac{(q+1)(q-4)}{12q},\qquad
 C_q=\frac{(q+1)(2q+1)}{12q}.
\end{equation}
The ratio is legitimate: numerator and denominator both start in degree
one and the denominator has nonzero coefficient $-q$ there.

\subsection{An all-orders factorization}

\begin{theorem}[Factorization and finite-cover obstruction for the model]
\label{thm:modelfactor}
The rescaled coordinate has the all-orders form
\begin{equation}\label{eq:modelfactor}
 \Psi_\eps(w)=w(1+w^q)^{\rho_q(\eps)}\ExpH(P_\eps(w)),
 \quad P_\eps\in\eps^2\C[w][[\eps]],\quad P_\eps(0)=0.
\end{equation}
At every nonzero fixed point the monodromy multiplier $M_q$ satisfies
\begin{equation}\label{eq:Mq}
 M_q=\ExpH(2\pi i\rho_q),\qquad
 M_q^q=1+\pi i(q+1)\eps+O(\eps^2).
\end{equation}
It has infinite order.  The leading coordinate becomes single-valued after
finite ramification, but no finite branched cover makes $\Psi_\eps$
single-valued at these points.
\end{theorem}

\begin{proof}
Apply Corollary~\ref{cor:factorization}.  All nonzero-root exponents coincide,
and
$\prod_{a^q=-1}(1-w/a)=1+w^q$.  The branches of the logarithms agree near
$0$, where each is normalized to vanish, so their sum is
$\log(1+w^q)$.  This gives~\eqref{eq:modelfactor} with the same polynomial
part $P_\eps$.
The coefficient at zero in $q\rho_q$ is $-1$ and the coefficient at $\eps$
is $(q+1)/2$.  Exponentiation proves~\eqref{eq:Mq}.
In fact no positive power is $1$, by Lemma~\ref{lem:torsion}, since
$A_q\ne0$.  The leading exponent $-1/q$ is rational; taking a $q$fold local
ramification removes its branching.  Theorem~\ref{thm:monodromy} proves
that the positive part of the full exponent survives every finite
ramification.
\end{proof}

A circuit repeated $q$ times thus returns the leading branch to itself but
still changes the full answer by a nonzero infinitesimal amount.  The
monodromy group records information that is lost under passage to standard
part.

\subsection{The first two correction functions}

Let $L(w)=\log(1+w^q)$ on a branch normalized by $L(0)=0$.
The next formula makes the logarithmic obstruction visible without any
operator notation.

\begin{proposition}[Explicit second-order expansion]
\label{prop:second}
On any admissible common domain,
\begin{equation}\label{eq:secondorder}
 \begin{split}
 \Psi_\eps(w)=w(1+w^q)^{-1/q}\biggl[1+A_q\eps L(w)
  +\eps^2\biggl(B_qL(w)+\frac{A_q^2}{2}L(w)^2-C_qw^q\biggr)
  +O(\eps^3)\biggr].
 \end{split}
\end{equation}
Equivalently, the polynomial term in~\eqref{eq:modelfactor} begins
$P_\eps(w)=-C_q\eps^2 w^q+O(\eps^3)$.
\end{proposition}

\begin{proof}
With $V=V_q$ and $\mu=1$, expansion of~\eqref{eq:alphaOmega} using
\eqref{eq:modifiedvector} gives
\[
 [\eps^2]\Omega_\eps
  =\frac{1/3-V'/4-(V')^2/12-VV''/12}{V}.
\]
Substitute $V=w(1+w^q)$ and collect powers of $w^q$.  The result is
\begin{equation}\label{eq:omega2model}
 [\eps^2]\Omega_\eps
  =-\frac{(q+1)(2q+1)}{12}w^{q-1}
   +\frac{(q+1)(q-4)}{12}\frac{w^{q-1}}{1+w^q}.
\end{equation}
The first summand is the polynomial part, whose normalized primitive is
$-C_qw^q$.  The second has normalized primitive $B_qL(w)$.
The coefficient at $\eps$ integrates to $A_qL(w)$, and the coefficient at
zero gives $\Psi_0$.  Exponentiating these corrections yields
\eqref{eq:secondorder}.
\end{proof}

For $q=1$ the formula is particularly simple:
\begin{equation}\label{eq:q1second}
 \Psi_\eps(w)=\frac{w}{1+w}
 \left[1+\eps\log(1+w)
 +\frac{\eps^2}{2}\bigl(\log^2(1+w)-\log(1+w)-w\bigr)
 +O(\eps^3)\right].
\end{equation}
For $q=4$ the coefficient $B_q$ happens to vanish, but the first-order
logarithm and hence the infinite-order monodromy do not.  An accidental
cancellation at one higher order cannot repair the obstruction.

\subsection{How the three representations fit together}

There are three useful descriptions of the same inner-disk coordinate.
The formal Taylor series~\eqref{eq:originalh} supplies exact valuations and
the sharp disk.  The common-domain coefficient series
$\Psi_\eps=\sum\eps^n\Psi_n$ supplies continuation at the rescaled boundary.
The factorization~\eqref{eq:modelfactor} isolates monodromy in a single Hahn
constant, with a single-valued polynomial-coefficient factor left over.

These descriptions answer different questions.  To evaluate very close to
$0$, the recurrence is efficient and directly summable.  To determine
branching near a boundary fixed point, the exact ratio of logarithms avoids
computing the many individual Taylor coefficients.  To understand why
algebraicity fails, even when the leading function is algebraic, the first
correction or the monodromy orbit is enough.

\section{Exact computation and reproducibility}
\label{sec:computation}

The archive includes \texttt{code/verify.py}, written for Python~3.9 or
later with the standard library only.  It uses \texttt{fractions.Fraction}
for every coefficient; it makes no floating-point comparison and needs no
network service, computer-algebra package, or proprietary evaluator.
Its recorded result is \texttt{data/verification.json}.  A small table of
computed coefficients is stored in \texttt{data/coefficients.csv}.

\subsection{What was checked}

The recorded run completed \textbf{2,764 exact assertions}.  The principal
ranges are shown in Table~\ref{tab:checks}.  ``Through $\eps^k$'' means that
all coefficients from degree zero through degree $k$ were compared exactly,
with the leading divisions handled before truncation.

\begin{table}[htbp]
\centering
\small
\begin{tabular}{@{}p{0.49\textwidth}p{0.43\textwidth}@{}}
\toprule
Identity or construction & Range of the recorded check\\
\midrule
Koenigs recurrence and leading binomial coefficients
 & $q=1,\ldots,8$; $m\le24$; $b_m$ through $\eps^8$\\[4pt]
Direct Schr\"oder residual, computed from the coefficients
 & Same $q,m$; residual through $\eps^9$\\[4pt]
First and second coefficient functions
 & Taylor coefficients through $w^{24q+1}$ for $q=1,\ldots,8$\\[4pt]
The logarithm ratio $\rho_q$
 & Constant, first and second coefficients for $q=1,\ldots,8$\\[4pt]
$\log C_g(p)=(\log C_g(w))p'$
 & $q=1,\ldots,6$; $p=w^n$, $0\le n\le6$; through $\eps^6$\\[4pt]
Generator coefficients and derivatives at fixed points
 & Modified field through $\eps^2$; fixed-point logarithms through $\eps^6$\\
\bottomrule
\end{tabular}
\caption{Finite exact consistency checks.  These are not a machine-checked
proof of the infinite statements.}
\label{tab:checks}
\end{table}

The recurrence and the directly substituted Schr\"oder residual are both
checked; agreement with a displayed formula is not the sole test.  The
vector-field calculation is separately performed using the operator
logarithm on sparse bivariate polynomials in $\eps$ and $w$.
For the nonzero model fixed points, derivatives are reduced using
$w^q=-1$, avoiding approximate complex root calculations.

\subsection{A reproducible finite algorithm}

To compute $b_0,\ldots,b_M$ modulo $\eps^{N+1}$, keep each coefficient as an
array of $N+1$ rational numbers.  In~\eqref{eq:brecurrence}, cancel the known
factor $\eps$ before division.  The remaining denominator has constant
coefficient $qm\ne0$, so triangular formal division is exact.
The numerator uses only already-computed $b_j$.  Powers of $1+\eps$ are
obtained from ordinary binomial coefficients.  Straight polynomial-array
arithmetic gives a simple finite algorithm; no assertion of optimal
complexity is needed for the reported ranges.

To compute the generator through $\eps^N$, form $N^r w$ successively and
stop at $r=N$.  Each application of $N=C_g-I$ increases the $\eps$ order by
at least one.  Truncation is therefore mathematically certified, rather
than a heuristic cutoff.  The combination
$\sum_{r=1}^N(-1)^{r+1}N^r w/r$ gives the logarithm to that order.

Running
\begin{verbatim}
python3 code/verify.py
\end{verbatim}
recreates the JSON report and the CSV table.  The supplied report records
Python~3.13.5.  A discrepancy is reported by a failing assertion rather than
silently rounded away.

\paragraph{What the program does not establish.}
No finite test proves Neumann summability, existence on an arbitrary
ordinary domain, compatibility for arbitrary value groups, or the
algebraicity dichotomy.  Those depend on the support arguments, exact
logarithmic residue identity, and field-automorphism proofs above.
The program also does not implement a complete symbolic surcomplex CAS or
a numerical evaluator at arbitrary surreal inputs.

\section{Scope, limitations, and further mathematical questions}
\label{sec:limits}

\subsection{The conclusions that have actually been obtained}

The main result is a uniform negative answer to an algebraicity question
for an entire polynomial family, not an isolated expansion.  Whenever
$V$ is nonlinear, the normalized coordinate of $w+\eps V(w)$ is
transcendental even over $\C(w)((t^\Gamma))$.  In the simple-root case the
obstruction has an exact local expression, and if the leading answer is
algebraic, the first correction is already not algebraic.
The common-domain theorem also applies to a controlled class of positive
Hahn perturbations beyond that polynomial family.  The model calculations
add an exact summability boundary and an isometric conjugacy on the inner
disk.

The most important hypotheses are not cosmetic.  The simple zero at the
normalization point permits the regular quotient in
\eqref{eq:regularR}.  A common ordinary domain permits coefficientwise
primitives without an uncontrolled sequence of shrinking radii.  Positive
support makes the operator logarithm strongly meaningful.  The polynomial
Euler form fixes all zeros and gives the exact factor $W_\eps=V B_\eps$;
a general perturbation can move zeros or create different coefficient pole
patterns.  Finally, divisibility of $\Gamma$ is used to keep the resonant
rescaling inside the chosen workspace.

\subsection{What remains outside the theorem}

This paper does not classify arbitrary surcomplex dynamical systems,
construct a global Julia or Fatou theory on the proper class $\No[i]$, or
identify the formal logarithm with a convergent autonomous flow at every
ordinary parameter.  It does not use ordinary topological convergence in
the full surcomplex plane.  Ordinary analytic continuation occurs only in
the spatial coefficient functions, and strong Hahn evaluation is a
separate, explicitly certified operation.

The transcendence result is specific to the indicated rational-coefficient
Hahn base.  It does not say that $H_\eps$ is differentially transcendental:
in fact it satisfies the first-order linear equation
$H_\eps'=\Omega_\eps H_\eps$, with $\Omega_\eps\in\B$.  Thus the construction
is differentially algebraic in a very explicit way while being
algebraically transcendental.  This distinction prevents an unwarranted
strengthening of the theorem.

\subsection{Directions suggested by the proofs}

A natural next problem is to characterize the monodromy of
$w+\eps V(w)+\eps^2R(w)$ when its fixed points move.  An analogue of the local derivative
identity at moving formal fixed points is a natural intermediate target;
translating it into a common-domain continuation theorem would also require
controlling those moving singularities.  The present Euler family
avoids that additional difficulty rather than concealing it.

In several variables, the logarithm-of-substitution construction still
suggests a vector field, but the scalar quotient $\alpha/W$ has no direct
replacement.  Resonant homological equations and noncommuting monodromy
would replace the one-dimensional integration argument.  A useful concrete
goal would be a support-controlled normal form for a diagonal multiplier
with infinitesimal resonances, with each obstruction assigned a certified
Hahn support.

For formalization in the repository, the finite recurrence, its leading
coefficient formula, and the exact valuation calculation are natural
algebraic entry points.  The next layer would formalize the common-domain
coefficient ring and its strong Taylor operator.  Ordinary covering-space
monodromy and the infinite-orbit algebraicity test should be kept as a
separate analytic layer, not hidden inside an unverified symbolic identity.
No Lean implementation of these new theorems is included or claimed here.

\paragraph{Conclusion.}
The leading complex coordinate and the full Hahn coordinate can agree on
a simple algebraic picture while disagreeing completely on finite
ramification.  The exact logarithm ratio at a fixed point is the quantity
that detects the difference.  It simultaneously explains the hidden
branching in the resonant model and proves the polynomial
Hahn-algebraicity dichotomy.

\appendix
\section{An elementary audit of the support arguments}
\label{app:support}

The following four checks summarize the support obligations in the proofs.
They are included to make it possible to audit every infinite operation
without treating ``formal'' as a blanket justification.

\paragraph{Substitution.}
For a fixed input support $T$ and positive displacement support $S$, a
Taylor term has exponent $\tau+s_1+\cdots+s_n$.  Neumann's lemma says that
at a fixed target exponent there are only finitely many choices of
$\tau$, $n$, and the word $(s_1,\ldots,s_n)$.  Differentiating the ordinary
coefficient does not alter the exponent.  Thus coefficient extraction
reduces substitution to finite algebra.

\paragraph{Operator logarithms.}
In $N^r\Phi$, at least $r$ positive displacement factors occur.  The same
finite-word certificate bounds $r$ at each target exponent.  Therefore
$\log(I+N)$, $(I+N)^{-1}$, and the binomial family $(I+N)^s$ are all
coefficientwise finite.  For the binomial proof of the derivation law, the
finite bound is independent of the ordinary parameter $s$, so coefficient
extraction genuinely gives a polynomial in $s$.

\paragraph{Normalized integration and exponentiation.}
Integration is performed on the fixed simply connected ordinary domain and
does not create Hahn exponents.  The leading coefficient is exponentiated
as one ordinary holomorphic function.  Only the positive remainder is
expanded as a Hahn exponential.  Confusing these two exponentiations would
incorrectly require an infinite family of exponent-zero terms to be
strongly summable.

\paragraph{Evaluation and reversion.}
An infinitesimal spatial input contributes positive letters from its own
support.  In a double series with nonnegative $\eps$ powers, the union of
those letters with $\{\delta\}$ gives a single positive well-ordered
control set.  Spatial degree bounds the number of such letters from below.
This makes both evaluation and formal reversion legitimate on the inner
disk.  At the boundary, that positivity in the rescaled spatial variable
is lost, and the same argument correctly ceases to apply.

All of these are set-sized constructions.  A larger surreal scale may
exist outside the workspace, but it is never needed to justify a sum, nor
is a proper-class family of coefficients ever formed.

\section{Notation and dependency map}

\begin{center}
\small
\begin{tabular}{@{}p{0.25\textwidth}p{0.67\textwidth}@{}}
\toprule
Notation & Meaning\\
\midrule
$\K=\C((t^\Gamma))$ & Constant Hahn field embedded in $\No[i]$.\\[3pt]
$\A(U)$ & Hahn series with all coefficients holomorphic on the same $U$.\\[3pt]
$\B=\C(w)((t^\Gamma))$ & Hahn field with rational-function coefficients; the algebraicity base.\\[3pt]
$\mathscr L$ & Meromorphic-coefficient Hahn field on the ordinary universal cover.\\[3pt]
$C_F$, $N$, $D$ & Taylor substitution, its difference from the identity, and its strong logarithm.\\[3pt]
$W_\eps$, $\alpha_\eps$ & Generator divided by $\eps$, and $\Log(1+\mu\eps)/\eps$.\\[3pt]
$\Omega_\eps=\alpha_\eps/W_\eps$ & Rational-coefficient logarithmic derivative of $H_\eps$.\\[3pt]
$\rho_a$, $M_a$ & Exact residue and coefficientwise monodromy multiplier at a simple zero.\\[3pt]
$h$, $\Psi_\eps$ & Original formal coordinate and rescaled common-domain coordinate of the model.\\
\bottomrule
\end{tabular}
\end{center}

The common-domain theorem depends on Neumann's lemma, the derivation
identity, and normalized ordinary integration.  The principal dichotomy
then uses polynomial divisibility of the generator, the fixed-point
logarithm identity, and either the elementary growth bound for algebraic
functions or the infinite-monodromy-orbit argument.  The sharp disk theorem
instead uses the coefficient recurrence and strong evaluation of formal
reversion.  These two proof branches meet in the resonant model but do not
logically substitute for one another.

\clearpage
\begin{thebibliography}{9}

\bibitem{Repo}
Vladimir Reshetnikov,
\emph{Surreal: surreal and surcomplex research reports and formalization},
GitHub repository, documentation map, manifest, and surcomplex analysis
source inspected at commit
\texttt{39f2be6667ade51bca2b45daa47e289d69c09764}, 21 September 2026.
\href{https://github.com/VladimirReshetnikov/Surreal/tree/39f2be6667ade51bca2b45daa47e289d69c09764/docs}{Frozen documentation tree}.
The reports are research drafts, not independent peer-reviewed validation.

\bibitem{Neumann}
B.~H. Neumann,
\emph{On ordered division rings},
Transactions of the American Mathematical Society \textbf{66} (1949),
202--252.
\href{https://doi.org/10.1090/S0002-9947-1949-0032593-5}{doi:10.1090/S0002-9947-1949-0032593-5}.

\bibitem{Lindahl}
Karl-Olof Lindahl,
\emph{Linearization in ultrametric dynamics in fields of characteristic
zero---equal characteristic case},
p-Adic Numbers, Ultrametric Analysis and Applications \textbf{1}(4)
(2009), 307--316.
\href{https://arxiv.org/abs/1111.1993}{arXiv:1111.1993};
\href{https://doi.org/10.1134/S2070046609040049}{doi:10.1134/S2070046609040049}.

\bibitem{GV}
Vassili Gelfreich and Arturo Vieiro,
\emph{Interpolating vector fields for near identity maps and averaging},
2017 preprint.
\href{https://arxiv.org/abs/1711.01983}{arXiv:1711.01983}.
Published version:
\href{https://doi.org/10.1088/1361-6544/aacb8e}{doi:10.1088/1361-6544/aacb8e}.

\bibitem{KKS}
Elliot Kaplan, Lothar Sebastian Krapp, and Michele Serra,
\emph{Decomposing the automorphism group of the surreal numbers},
\href{https://arxiv.org/abs/2509.22374v3}{arXiv:2509.22374v3},
revised 23 April 2026.

\end{thebibliography}
\end{document}
